\chapter{Literature Review and Theoretical Foundations}
\label{ch:funda}

Detecting organized financial fraud on credit and transaction networks is fundamentally a graph problem. Fraudsters rarely operate in isolation; they coordinate across accounts, share devices, and intentionally interact with legitimate users to blend in. In this chapter, we examine how graph neural approaches have evolved to tackle this challenge, moving from early camouflage-filtering models (2020--2021) to recent spectral and temporal frameworks (2022--2026). We then review ant colony optimization and stigmergy, showing where neural-ACO hybrids currently succeed in combinatorial routing and why integrating persistent pheromone memory into GNN message-passing addresses a major blind spot in financial fraud detection.

\section[GNN-Based Fraud Detection]{\textbf{GNN-Based Fraud Detection (Survey 2020--2026)}}
\label{sec:gnn_fraud_survey}

\subsection{Foundational Fraud GNN Architectures (2020--2021)}
When researchers first applied standard Graph Convolutional Networks (GCN) \cite{kipf2017semi} and GraphSAGE \cite{hamilton2017inductive} to transaction graphs, they quickly ran into a core limitation. Standard message-passing relies on the homophily assumption---that connected nodes share similar attributes and labels. In fraud networks, the opposite is true: fraudsters deliberately connect to benign accounts (camouflage) and conceal their similarities. Uniform neighborhood averaging ends up smoothing away anomaly signals, pulling fraud embeddings into the benign majority cluster. To fix this, several specialized architectures emerged between 2020 and 2021:

\begin{enumerate}
    \item \textbf{CARE-GNN (Dou et al., KDD 2020 \cite{dou2020enhancing}):} Dou et al. were the first to tackle camouflage directly on multi-relational graphs. Their key insight was to pair relation-aware attention with a Reinforcement Learning (RL) agent. Instead of aggregating every neighbor, the RL controller adjusts a similarity threshold per relation type based on feature similarity $\mathcal{D}(h_u, h_v)$:
    \begin{equation}
        p^{(r, l)} \leftarrow p^{(r, l)} + \Delta p^{(r, l)} \quad \text{via RL policy gradient}
    \end{equation}
    While effective at filtering obvious feature outliers, CARE-GNN makes instantaneous, node-by-node decisions each epoch. It maintains no persistent trail memory across training steps, meaning multi-hop structural evidence cannot accumulate over time.
    \item \textbf{PC-GNN (Liu et al., WWW 2021 \cite{liu2021pick}):} Real-world fraud datasets typically contain fewer than 10\% illicit nodes. To prevent models from collapsing into majority-class predictions, Liu et al. introduced a label-balanced ``Pick and Choose'' sampling strategy. PC-GNN oversamples fraud nodes while filtering out uninformative benign neighbors during mini-batch generation. However, it relies on heuristic feature similarity, which struggles when fraudsters forge realistic user attributes.
    \item \textbf{BWGNN (Tang et al., ICDM 2021 \cite{tang2021rethinking}):} Tang et al. analyzed fraud graphs from a spectral perspective. In homophilic graphs, the graph Laplacian is dominated by low-frequency components. In contrast, fraud edges introduce sharp feature discrepancies that manifest as high-frequency spectral signals. BWGNN captures these localized anomalies using a spectral Beta Wavelet filter:
    \begin{equation}
        h_i^{(l+1)} = \sum_{k=0}^K \beta_k \cdot \mathbf{L}^k h_i^{(l)}
    \end{equation}
    While effective for localized heterophilic transitions, BWGNN assumes a static graph structure and lacks an explicit mechanism to trace multi-hop fraud rings.
    \item \textbf{FRAUDRE (Zhang et al., ICDM 2021 \cite{zhang2021fraudre}):} Fraudsters often masquerade as high-degree benign hubs to dilute detection signals. Zhang et al. proposed FRAUDRE to penalize nodes that exploit high connectivity across relations, combined with a cost-sensitive reweighting loss to balance training.
\end{enumerate}

\subsection{Modern Fraud GNN Architectures (2022--2026)}
Between 2022 and 2026, research shifted toward handling low homophily directly, modeling group interactions, and handling noisy dynamic streams:

\begin{enumerate}
    \item \textbf{H2-FDetector (Shi et al., WWW 2022 \cite{shi2022h2}):} Shi et al. addressed the heterophily problem by separating neighborhood aggregation into two distinct streams: label-consistent (homophilic) and label-divergent (heterophilic) connections:
    \begin{equation}
        h_i = \text{CONCAT}\left(h_i^{\text{homo}}, h_i^{\text{hetero}}\right)
    \end{equation}
    \item \textbf{GAGA (Wang et al., WWW 2023 \cite{wang2023gaga}):} Moving beyond pairwise edge comparisons, Wang et al. designed a group aggregation transformer to detect synchronized collective behaviors across multiple accounts.
    \item \textbf{MTPNet (Zhang et al., IEEE TIFS, Sept 2025 \cite{mtpnet2025}):} Real financial transactions arrive as continuous streams. Zhang et al. partitioned transaction logs into multi-scale time windows and applied partitioned temporal attention. Yet attention scores are still recomputed independently per forward pass without a cumulative trail.
    \item \textbf{THG-OAFN (Wei \& Lee, Dec 2025 \cite{wei2025internet}):} Combined temporal heterogeneous graphs with a GraphSMOTE oversampling algorithm on YelpChi. However, generating synthetic nodes in financial credit networks can introduce topological artifacts that do not reflect actual borrower behavior.
    \item \textbf{Adversarial Generative Defense (Wu et al., IEEE TKDE, Feb 2026 \cite{wu2026defending}):} Wu et al. used a generative graph autoencoder to model adversarial edge additions. While it improved robustness by 3\%--7\%, generating and evaluating candidate edges scales quadratically as $\mathcal{O}(N^2)$, making it impractical for multi-million edge graphs.
    \item \textbf{Theoretical Noise Analysis (Ju et al., IEEE TPAMI, March 2026 \cite{ju2026survey}):} Ju et al. proved mathematically that standard GNN aggregation degrades exponentially once camouflage noise exceeds 20\%, reinforcing the need for edge-weight regularization that persists across training iterations.
\end{enumerate}

\subsection{Base Review Paper: Cheng et al. (Frontiers of Computer Science 2025)}
Our work draws direct guidance from the comprehensive survey by \textbf{Cheng et al. (2025)} \cite{cheng2025graph}, titled \textit{``Graph Neural Networks for Financial Fraud Detection: A Review''}, published in \textit{Frontiers of Computer Science}. Reviewing over 100 fraud detection architectures across retail credit, e-commerce, and anti-money laundering, the authors highlighted three critical takeaways:
\begin{enumerate}
    \item \textbf{Camouflage Degrades Performance Severely:} When fraud syndicates inject 20\%--40\% benign connections, the Macro-F1 of top-tier models (including CARE-GNN and PC-GNN) drops sharply by 25\%--42\%.
    \item \textbf{Standardized Multi-Relational Benchmarks:} Realistic evaluation requires typed relational networks. Benchmarks such as \textbf{YelpChi} (review graph with R-U-R, R-T-R, and R-S-R meta-paths) and \textbf{Amazon} (with U-P-U, U-S-U, and U-V-U relations) represent the standard testbeds.
    \item \textbf{The Need for Persistent, Bio-Inspired Priors:} Cheng et al. noted that pure gradient-based backpropagation easily settles into local camouflage traps. They advocated exploring \textit{adaptive bio-inspired heuristics and persistent structural prior mechanisms} capable of dynamically re-weighting edges across training cycles without incurring quadratic computational overhead. The ACO-GNN framework was designed directly in response to this recommendation.
\end{enumerate}

% ------------------------------------------------------------------------------
% SECTION 2.2: ANT COLONY OPTIMIZATION IN GRAPH LEARNING
% ------------------------------------------------------------------------------
\section[Ant Colony Optimization in Graph Learning]{\textbf{Ant Colony Optimization and Stigmergic Learning}}
\label{sec:aco_background}

Ant Colony Optimization (ACO) originates from Marco Dorigo's 1992 work \cite{dorigo1992optimization, dorigo1996ant} studying how biological ant colonies locate optimal foraging routes without a central controller. The underlying mechanism is *stigmergy*: indirect coordination mediated through physical modifications of the environment. When an ant finds food, it marks its traversal path with a volatile chemical secretion (pheromone). Subsequent ants detect this marker and tend to follow stronger scent trails over unvisited terrain. Over repeated traversals, shorter and more productive paths accumulate denser pheromone deposits, while longer or erroneous excursions fade due to natural evaporation. Stigmergy essentially acts as a decentralized, persistent spatial memory.

\subsection{Foundational Metaheuristic Dynamics}
In canonical Ant System (AS) and the Max-Min Ant System (MMAS) \cite{stutzle2000max}, an artificial ant $k$ situated at node $i$ chooses its next step $j \in \mathcal{N}(i)$ using a stochastic transition rule that balances historical trail intensity against immediate heuristic desirability:
\begin{equation}
    P_{ij}^{(k)} = \frac{\left[\tau_{ij}\right]^\alpha \cdot \left[\eta_{ij}\right]^\beta}{\sum_{l \in \mathcal{N}(i)} \left[\tau_{il}\right]^\alpha \cdot \left[\eta_{il}\right]^\beta}
    \label{eq:canonical_aco_transition}
\end{equation}
where $\tau_{ij}$ represents the pheromone level on edge $(i, j)$, and $\eta_{ij}$ provides problem-specific heuristic guidance (such as inverse distance $1/d_{ij}$ in routing). The exponents $\alpha, \beta \ge 0$ tune the balance between collective historical experience and greedy local cues.

Without continuous decay, early exploration choices would permanently dominate the network. To maintain adaptability and prevent premature stagnation, trails undergo exponential evaporation at a rate $\rho \in (0, 1)$:
\begin{equation}
    \tau_{ij} \leftarrow (1 - \rho) \cdot \tau_{ij} + \sum_{k=1}^{M} \Delta \tau_{ij}^{(k)}
    \label{eq:canonical_evaporation}
\end{equation}
where $\Delta \tau_{ij}^{(k)}$ denotes the pheromone increment left by ant $k$, typically scaled by solution quality $Q / L_k$. In MMAS, explicit bounds $[\tau_{\min}, \tau_{\max}]$ are placed on every trail. This avoids two failure modes: trail starvation (where an edge's score drops to zero and becomes permanently unreachable) and trail saturation (where an early suboptimal path locks out all alternatives).

\begin{figure}[!htbp]
\centering
\resizebox{0.92\textwidth}{!}{%
\begin{tikzpicture}[
    font=\small\sffamily,
    fraud/.style={circle, draw=red!80, fill=red!15, very thick, minimum size=1.05cm, align=center, font=\bfseries\small},
    benign/.style={circle, draw=blue!75, fill=blue!12, very thick, minimum size=1.05cm, align=center, font=\bfseries\small},
    rfraud/.style={circle, draw=red!90, fill=red!25, very thick, minimum size=1.05cm, align=center, font=\bfseries\small},
    rbenign/.style={circle, draw=blue!80, fill=blue!10, very thick, minimum size=1.05cm, align=center, font=\bfseries\small},
    arr_norm/.style={draw=gray!70, -{Latex[length=2.2mm]}, thick},
    arr_camo/.style={draw=orange!80!black, -{Latex[length=2.2mm]}, dashed, thick},
    arr_strong/.style={draw=red!85!black, -{Latex[length=2.8mm]}, line width=2.4pt},
    arr_faded/.style={draw=gray!35, -{Latex[length=1.8mm]}, loosely dashed, line width=0.7pt},
    box_panel/.style={draw=gray!40, rounded corners=8pt, fill=gray!2, inner sep=10pt}
]

%% ==================== LEFT PANEL: INITIAL STATE (t=0) ====================
\node[fraud] (L_F1) at (0.5, 2.2) {$F_1$};
\node[fraud] (L_F2) at (2.7, 2.2) {$F_2$};
\node[fraud] (L_F3) at (2.7, 0) {$F_3$};
\node[fraud] (L_F4) at (0.5, 0) {$F_4$};

\node[benign] (L_B1) at (-1.5, 1.1) {$B_1$};
\node[benign] (L_B2) at (4.5, 1.1) {$B_2$};

% Fraud ring edges (initial uniform)
\draw[arr_norm] (L_F1) -- node[above, font=\scriptsize]{$\tau_0{=}1.0$} (L_F2);
\draw[arr_norm] (L_F2) -- node[right, font=\scriptsize]{$\tau_0{=}1.0$} (L_F3);
\draw[arr_norm] (L_F3) -- node[below, font=\scriptsize]{$\tau_0{=}1.0$} (L_F4);
\draw[arr_norm] (L_F4) -- node[left, font=\scriptsize]{$\tau_0{=}1.0$} (L_F1);

% Camouflage edges
\draw[arr_camo] (L_F1) -- node[above left, font=\scriptsize, text=orange!80!black]{Camo} (L_B1);
\draw[arr_camo] (L_F4) -- node[below left, font=\scriptsize, text=orange!80!black]{Camo} (L_B1);
\draw[arr_camo] (L_F2) -- node[above right, font=\scriptsize, text=orange!80!black]{Camo} (L_B2);
\draw[arr_camo] (L_F3) -- node[below right, font=\scriptsize, text=orange!80!black]{Camo} (L_B2);

% Label for left panel
\node[below=0.7cm of L_F4, xshift=1.1cm, align=center, font=\bfseries\small] (lbl_left) 
    {(a) Initial Transaction Topology ($t=0$)\\[2pt]
     \normalfont\scriptsize Camouflage edges disguise collusive ring;\\
     uniform initial priors $\tau_{ij} = \tau_0 = 1.0$};

\begin{scope}[on background layer]
    \node[box_panel, fit=(L_B1)(L_B2)(L_F1)(L_F3)(lbl_left)] (panel_left) {};
\end{scope}

%% ==================== CENTER TRANSITION ARROW ====================
\draw[-{Latex[length=4mm, width=3mm]}, line width=3pt, draw=teal!80] 
    (5.4, 1.1) -- (7.6, 1.1)
    node[midway, above=3pt, align=center, font=\scriptsize\bfseries, text=teal!90!black] 
    {Iterative Stigmergy\\[1pt]
     $\tau \leftarrow (1{-}\rho)\tau + \Delta\tau$};

\node[align=center, font=\tiny, text=gray!70, below=3pt] at (6.5, 0.8) 
    {Reinforce suspicious edges\\Evaporate camouflage};

%% ==================== RIGHT PANEL: CONVERGED STATE (t=T) ====================
\node[rfraud] (R_F1) at (10.2, 2.2) {$F_1$};
\node[rfraud] (R_F2) at (12.4, 2.2) {$F_2$};
\node[rfraud] (R_F3) at (12.4, 0) {$F_3$};
\node[rfraud] (R_F4) at (10.2, 0) {$F_4$};

\node[rbenign] (R_B1) at (8.2, 1.1) {$B_1$};
\node[rbenign] (R_B2) at (14.2, 1.1) {$B_2$};

% Fraud ring edges (strongly reinforced)
\draw[arr_strong] (R_F1) -- node[above, font=\scriptsize\bfseries, text=red!80!black]{$\tau \to \tau_{\max}$} (R_F2);
\draw[arr_strong] (R_F2) -- node[right, font=\scriptsize\bfseries, text=red!80!black]{$\tau \to \tau_{\max}$} (R_F3);
\draw[arr_strong] (R_F3) -- node[below, font=\scriptsize\bfseries, text=red!80!black]{$\tau \to \tau_{\max}$} (R_F4);
\draw[arr_strong] (R_F4) -- node[left, font=\scriptsize\bfseries, text=red!80!black]{$\tau \to \tau_{\max}$} (R_F1);

% Camouflage edges (evaporated)
\draw[arr_faded] (R_F1) -- node[above left, font=\tiny, text=gray!60]{$\tau \to \tau_{\min}$} (R_B1);
\draw[arr_faded] (R_F4) -- node[below left, font=\tiny, text=gray!60]{$\tau \to \tau_{\min}$} (R_B1);
\draw[arr_faded] (R_F2) -- node[above right, font=\tiny, text=gray!60]{$\tau \to \tau_{\min}$} (R_B2);
\draw[arr_faded] (R_F3) -- node[below right, font=\tiny, text=gray!60]{$\tau \to \tau_{\min}$} (R_B2);

% Label for right panel
\node[below=0.7cm of R_F4, xshift=1.1cm, align=center, font=\bfseries\small] (lbl_right) 
    {(b) Converged Stigmergic Topology ($t=T$)\\[2pt]
     \normalfont\scriptsize Multi-hop fraud ring reinforced ($\tau_{ij} \uparrow$);\\
     camouflage pathways dissolved ($(1{-}\rho)\tau \to \tau_{\min}$)};

\begin{scope}[on background layer]
    \node[box_panel, fit=(R_B1)(R_B2)(R_F1)(R_F3)(lbl_right)] (panel_right) {};
\end{scope}

\end{tikzpicture}%
}
\caption[Conceptual Stigmergic Disentanglement of Collusive Fraud Rings]{Conceptual mechanism of biological stigmergy and pheromone evaporation under deceptive camouflage. (a) In the raw transaction graph, collusive fraudsters ($F_1$--$F_4$) disguise their connectivity by injecting camouflage edges to benign accounts ($B_1, B_2$). (b) Over iterative training epochs, correlated fraud predictions trigger pheromone reinforcement along ring edges ($\tau \to \tau_{\max}$), whereas non-reinforced camouflage pathways continuously evaporate ($(1-\rho)\tau \to \tau_{\min}$), isolating the collusive syndicate without heuristic edge pruning.}
\label{fig:stigmergy_camouflage_concept}
\end{figure}

Figure~\ref{fig:stigmergy_camouflage_concept} illustrates how this continuous reinforcement and evaporation mechanism is adapted to counter adversarial camouflage in financial transaction networks.

\subsection{Neural Combinatorial Optimization Hybrids}
Over the past five years, several pioneering works have combined deep representation learning with ant systems:
\begin{enumerate}
    \item \textbf{DeepACO (NeurIPS 2023)} \cite{ye2023deepaco}: Targets NP-hard routing problems like the TSP, CVRP, and Orienteering. Instead of computing static geometric heuristics $\eta(i, j)$, DeepACO uses a GNN to generate learned heuristic matrices from node embeddings, which artificial ants then sample to construct tours.
    \item \textbf{GFACS (ICLR 2024)} \cite{gfacs2024}: Incorporates Graph Flow Attention into Ant Colony Systems, dynamically parameterizing construction graphs to solve constrained combinatorial problems under operational constraints.
    \item \textbf{NeuFACO (2025)} \cite{neufaco2025}: Introduces neural-guided fast ACO algorithms for dynamic vehicle routing with real-time demand fluctuations.
\end{enumerate}

\subsection{The Unexplored Research Gap in Financial Fraud Detection}
Notice an important distinction: all of the above frameworks focus strictly on \textit{discrete combinatorial optimization} (ordering paths or selecting vertex subsets). This leaves an obvious research gap:
\begin{quote}\small\raggedright
    \textit{Prior to this work, no literature has integrated continuous stigmergic pheromone dynamics directly into Graph Neural Network neighborhood message-passing for graph anomaly detection or financial fraud ring identification.}
\end{quote}
Existing fraud detection graph architectures rely on three main strategies:
\begin{itemize}
    \item \textbf{Reinforcement Learning Filtering (e.g., CARE-GNN \cite{dou2020enhancing}):} Uses multi-agent RL (specifically Bernoulli multi-armed bandits) to select a neighbor filtering threshold per relation. While helpful for dropping local feature outliers, the RL agent makes greedy, single-step decisions and maintains no cross-epoch structural memory.
    \item \textbf{Spectral Wavelet / Band-pass Filtering (e.g., BWGNN \cite{tang2021rethinking}):} Designs localized spectral band-pass filters to isolate high-frequency feature shifts caused by camouflage. However, spectral approaches treat all graph edges uniformly without tracking persistent transaction patterns over progressive dynamic updates.
    \item \textbf{Transformer Group Aggregations (e.g., GAGA \cite{wang2023gaga}):} Employs all-to-all dense attention to bypass low homophily. This incurs quadratic $\mathcal{O}(N^2)$ computational complexity and discards edge-relational inductive biases.
\end{itemize}

None of the above frameworks maintain an \textbf{iterative persistent structural memory} that accumulates collective evidence over training epochs while steadily evaporating deceptive camouflage connections. By adapting stigmergic pheromone trails as an inductive structural prior alongside parametric neural attention, our framework directly bridges this gap.


% ------------------------------------------------------------------------------
% SECTION 2.3: RESEARCH GAP IDENTIFICATION
% ------------------------------------------------------------------------------
\section[Research Gap Identification]{\textbf{Research Gap Identification}}
\label{sec:research_gap}

To synthesize the contemporary landscape, Table~\ref{tab:lit_comparison} presents a systematic comparative analysis across 12 prominent fraud detection frameworks.

\begin{table}[htb]
\centering
\caption[Comparison of graph fraud detection frameworks against ACO-GNN]{Systematic comparison of graph fraud detection frameworks against ACO-GNN.}
\label{tab:lit_comparison}
\scriptsize
\setlength{\tabcolsep}{2.2pt}
\begin{tabularx}{\textwidth}{@{}l c c c c c >{\raggedright\arraybackslash}X@{}}
\toprule
\textbf{Model} & \textbf{Venue / Year} & \textbf{Graph Type} & \textbf{Camouflage} & \textbf{Memory} & \textbf{Bio-Inspired} & \textbf{Operational Bottleneck} \\
\midrule
GCN \cite{kipf2017semi} & ICLR 2017 & Homogeneous & None & None & No & Feature over-smoothing under structural noise \\
GAT \cite{velivckovic2018graph} & ICLR 2018 & Homogeneous & Feature Attention & None & No & Attention deceived by attribute camouflage \\
CARE-GNN \cite{dou2020enhancing} & KDD 2020 & Heterogeneous & RL Threshold & None & No & Greedy per-node policy, high RL instability \\
PC-GNN \cite{liu2021pick} & WWW 2021 & Heterogeneous & Subgraph Sampling & None & No & Relies on surface feature similarity \\
BWGNN \cite{tang2021rethinking} & ICDM 2021 & Homogeneous & Beta Wavelets & None & No & Static spectral filter; misses deep rings \\
FRAUDRE \cite{zhang2021fraudre} & ICDM 2021 & Heterogeneous & Cost Reweight & None & No & False alarms on high-degree benign hubs \\
H2-FDetector \cite{shi2022h2} & WWW 2022 & Heterogeneous & Disentangle & None & No & Static separation; no multi-hop persistence \\
GAGA \cite{wang2023gaga} & WWW 2023 & Heterogeneous & Group Attention & None & No & High computational overhead for group mining \\
DeepACO \cite{ye2023deepaco} & NeurIPS 2023 & Homogeneous & N/A & TSP Trails & Yes & Limited strictly to combinatorial routing \\
MTPNet \cite{mtpnet2025} & IEEE TIFS 2025 & Temporal & Partition Attn & None & No & Recomputed per-pass; lacks cross-epoch trail \\
THG-OAFN \cite{wei2025internet} & PLOS ONE 2025 & Heterogeneous & GraphSMOTE & None & No & Oversampling causes topological artifacts \\
Wu et al. \cite{wu2026defending} & IEEE TKDE 2026 & Heterogeneous & Generative AE & Latent & No & $\mathcal{O}(N^2)$ candidate edge generation cost \\
\midrule
\textbf{ACO-GNN} & \textbf{This Work} & \textbf{Heterogeneous} & \textbf{Pheromones} & \textbf{Stigmergic} & \textbf{Yes} & \textbf{Blends attention with stigmergy; $\mathcal{O}(|E|)$} \\
\bottomrule
\end{tabularx}
\end{table}

This taxonomy highlights three fundamental gaps that motivated our proposed architecture:
\begin{enumerate}
    \item \textbf{Absence of Persistent Edge Memory in GNN Aggregation:} Existing fraud GNNs treat message-passing as a memoryless process. Edge aggregation weights are computed in isolation at each iteration from instantaneous node embeddings $h_u, h_v$. When fraudsters disguise their profile features, the attention mechanism falls into the same camouflage traps in every epoch. The model lacks a cumulative trace that remembers persistent suspicious multi-hop connectivity.
    \item \textbf{Brittleness of Binary Edge Pruning:} Contemporary camouflage defenses (such as CARE-GNN's RL thresholding) rely on hard binary edge pruning ($e_{ij} \in \{0, 1\}$). Dropping edges prematurely often severs multi-hop pathways necessary to identify coordinated fraud syndicates spanning 3 to 6 hops.
    \item \textbf{Unexplored Stigmergic Priors in Deep Learning:} While ACO has proven remarkably effective for discrete combinatorial routing, its continuous stigmergic reinforcement and evaporation dynamics have never been integrated into the continuous message aggregation equations of Graph Neural Networks for node classification or fraud detection.
\end{enumerate}
